\documentclass[11pt]{article}

\usepackage[T1]{fontenc}
\usepackage[utf8]{inputenc}
\usepackage{lmodern}
\usepackage[margin=1.1in]{geometry}
\usepackage{amsmath,amssymb,amsthm,mathtools}
\usepackage{enumitem}
\usepackage{microtype}
\usepackage{xcolor}
\usepackage{titlesec}
\usepackage{hyperref}

\titleformat{\section}
  {\normalfont\LARGE\bfseries}{\thesection}{1em}{}
\titleformat{\subsection}
  {\normalfont\Large\bfseries}{\thesubsection}{1em}{}

\hypersetup{
  colorlinks=true,
  linkcolor=blue!55!black,
  citecolor=green!40!black,
  urlcolor=blue!60!black
}

\newtheorem{theorem}{Theorem}[section]
\newtheorem{proposition}[theorem]{Proposition}
\newtheorem{lemma}[theorem]{Lemma}
\newtheorem{corollary}[theorem]{Corollary}
\theoremstyle{remark}
\newtheorem*{remark}{Remark}
\theoremstyle{definition}
\newtheorem{definition}[theorem]{Definition}

\DeclareMathOperator{\Spec}{Spec}
\DeclareMathOperator{\Frac}{Frac}
\DeclareMathOperator{\Gal}{Gal}
\DeclareMathOperator{\Hom}{Hom}
\DeclareMathOperator{\Aut}{Aut}
\DeclareMathOperator{\Ram}{Ram}
\DeclareMathOperator{\Norm}{Norm}
\DeclareMathOperator{\Ann}{Ann}
\DeclareMathOperator{\core}{core}

\newcommand{\A}{\mathbb A}
\newcommand{\C}{\mathbb C}
\newcommand{\Obs}{\operatorname{Obs}}
\newcommand{\id}{\mathrm{id}}
\newcommand{\qedbox}{\hfill$\square$}
\newcommand{\leanpath}[1]{\texttt{\nolinkurl{#1}}}

\newenvironment{leanfiles}
  {\par\medskip\begingroup\footnotesize\raggedright\noindent
   \textit{Lean correspondence.}\ }
  {\par\endgroup\medskip}

\title{Collision Ideals and Off-Diagonal Sheets}
\author{Chloe van der Vlugt\\Codex 5.6 Sol}
\date{\today}

\begin{document}

\maketitle

\begin{abstract}
For a polynomial map
\(F=(F_1,\dots,F_n)\colon\A^n_{\C}\to\A^n_{\C}\), define in
\(S=\C[x_1,\dots,x_n,y_1,\dots,y_n]\) the collision and diagonal ideals
\[
  I_R=\bigl(F_i(x)-F_i(y)\bigr)_{i=1}^n,
  \qquad
  I_\Delta=(x_i-y_i)_{i=1}^n.
\]
The inclusion \(I_R\subseteq I_\Delta\) defines the obstruction module
\(I_\Delta/I_R\), and
\[
  I_\Delta/I_R=0
  \quad\Longleftrightarrow\quad
  I_R=I_\Delta
  \quad\Longleftrightarrow\quad
  F\text{ is a polynomial automorphism}.
\]
In dimension two, let \(\mathcal D_F\) be the finite-cover diagram of a
map with constant nonzero Jacobian determinant, and let
\(\mathcal H(\mathcal D_F)\) be its hidden-inertia locus.  Then
\[
  \mathcal H(\mathcal D_F)=\varnothing
  \quad\Longrightarrow\quad
  I_R=I_\Delta.
\]
For the separable nonnormal cubic class in dimension three, let
\[
  K=\C(F_1,F_2,F_3),
  \qquad
  L=\C(x_1,x_2,x_3),
\]
and let \(N\) be the normal closure of \(L/K\).  Then
\[
  L\otimes_K L\simeq L\times N,
  \qquad
  \Gal(N/K)\simeq S_3.
\]
The off-diagonal factor \(N\) is the ordered-root Galois-closure sheet.
Consequently, \(I_\Delta/I_R\neq0\).  These constructions and
implications are formalized in Lean~4, with the external inputs isolated
explicitly.
\end{abstract}

\begin{center}
  \small Licensed under
  \href{https://creativecommons.org/licenses/by/4.0/}
    {CC BY 4.0}.
\end{center}

\tableofcontents

\section{General construction}
\label{sec:general-construction}

\subsection{The two ideals and their obstruction}
\label{sec:collision-ideals}

Throughout the paper, the base field is \(\C\).  Let
\[
  F=(F_1,\dots,F_n)\colon
  X=\A^n_{\C}\longrightarrow\A^n_{\C}
\]
be a polynomial map, and let
\[
  S=\C[x_1,\dots,x_n,y_1,\dots,y_n].
\]
The map \(F\) is Keller if
\[
  \det JF\in\C^\times.
\]
Define
\[
  I_R(F)=\bigl(F_i(x)-F_i(y)\bigr)_{i=1}^{n},
  \qquad
  I_\Delta=(x_i-y_i)_{i=1}^{n},
  \qquad
  \Obs(F):=I_\Delta/I_R(F).
\]
For a fixed map \(F\), we abbreviate \(I_R(F)\) to \(I_R\).

\begin{proposition}[Canonical containment]
\label{prop:collision-contained-in-diagonal}
For every polynomial map \(F\),
\[
  I_R(F)\subseteq I_\Delta.
\]
\end{proposition}

\begin{proof}
For \(H\in\C[x_1,\dots,x_n]\), the difference \(H(x)-H(y)\) belongs to
\((x_1-y_1,\dots,x_n-y_n)\).  This follows by induction on \(H\), or by
telescoping each monomial one coordinate at a time.  Apply the statement
to each coordinate \(F_i\).
\end{proof}

Let \(C_F:=S/I_R\) and
\(A:=\C[x_1,\dots,x_n]\).  Diagonal evaluation induces a surjective ring
map
\[
  \bar\mu_F\colon C_F\longrightarrow A,
  \qquad
  [s]_{I_R}\longmapsto s(x,x).
\]

\begin{proposition}[The diagonal kernel]
\label{prop:obstruction-kernel}
As ideals of \(C_F\),
\[
  \ker(\bar\mu_F)
  =
  \Obs(F)
  =
  I_\Delta/I_R.
\]
Consequently,
\[
  \ker(\bar\mu_F)=0
  \quad\Longleftrightarrow\quad
  \Obs(F)=0
  \quad\Longleftrightarrow\quad
  I_R=I_\Delta.
\]
\end{proposition}

\begin{proof}
The kernel of diagonal evaluation \(S\to A\) is \(I_\Delta\).  Passing
through \(S/I_R\) therefore identifies the displayed kernel with
\(I_\Delta/I_R\).  Since \(I_R\subseteq I_\Delta\),
\[
  I_\Delta/I_R=0
  \quad\Longleftrightarrow\quad
  I_R=I_\Delta.
\]
\end{proof}

\subsection{The fiber-product interpretation}
\label{sec:fiber-product}

Let
\[
  A=\C[x_1,\dots,x_n],\qquad
  B=\C[F_1,\dots,F_n]\subseteq A.
\]
Let \(X=\Spec A\), let \(Y=\Spec B\), and let
\[
  f_F\colon X\longrightarrow Y
\]
be the morphism induced by \(B\hookrightarrow A\).
The homomorphisms
\(\iota_x,\iota_y\colon A\rightrightarrows C_F\), defined by
\(\iota_x(a)=[a(x)]\) and \(\iota_y(a)=[a(y)]\), agree on \(B\) because
\([F_i(x)]=[F_i(y)]\) in \(C_F\).  Their common restriction equips
\(C_F\) with a \(B\)-algebra structure.

\begin{theorem}[Fiber-product interpretation]
\label{thm:collision-tensor-product}
There is a canonical \(B\)-algebra equivalence and its dual
affine-scheme equivalence
\[
  C_F\cong A\otimes_B A,
  \qquad
  R_F:=\Spec C_F\cong X\times_YX.
\]
Under the first equivalence, \(\bar\mu_F\) is tensor multiplication
\[
  A\otimes_BA\longrightarrow A,\qquad a\otimes a'\longmapsto aa'.
\]
Its dual morphism is the diagonal
\[
  \Delta_{X/Y}\colon X\longrightarrow X\times_YX.
\]
\end{theorem}

\begin{proof}
A map \(C_F\to T\) to a commutative \(\C\)-algebra \(T\) is equivalent
to a pair of maps \(f,g\colon A\rightrightarrows T\) satisfying
\[
  f(F_i)=g(F_i)\qquad(1\leq i\leq n).
\]
Since the \(F_i\) generate \(B\), this is precisely the universal
property of \(A\otimes_BA\).  Applying \(\Spec\) to this equivalence and
using
\[
  \Spec(A\otimes_BA)
  \cong
  \Spec A\times_{\Spec B}\Spec A
\]
gives the asserted fiber-product equivalence.
The statements about multiplication and the diagonal follow on pure
tensors.
\end{proof}

\[
  R_F(\C)
  =
  \{(u,v)\in X(\C)\times X(\C):F(u)=F(v)\}.
\]
Inside
\[
  X\times X=\Spec S,
\]
the ideals \(I_R\) and \(I_\Delta\) define, respectively,
\(X\times_YX\) and the image of \(\Delta_{X/Y}\).  The
contravariance of \(I\mapsto V(I)\) therefore gives the equivalences
\[
  I_R\subseteq I_\Delta
  \quad\Longleftrightarrow\quad
  \Delta_{X/Y}(X)\subseteq X\times_YX
\]
and
\[
  I_R=I_\Delta
  \quad\Longleftrightarrow\quad
  X\times_YX=\Delta_{X/Y}(X)
\]
as closed subschemes of \(X\times X\).  Thus the algebraic and
geometric gaps are two descriptions of the same scheme-theoretic
containment.

\subsection{The off-diagonal factor and secant projectors}
\label{sec:secant-projector}

For every polynomial map \(F\), define
\[
  I_{\mathrm{off}}:=I_R:I_\Delta,
  \qquad
  C_F^\circ:=S/I_{\mathrm{off}},
  \qquad
  R_F^\circ:=\Spec C_F^\circ.
\]
This definition does not require the diagonal to be clopen.

\begin{corollary}[Off-diagonal emptiness criterion]
\label{cor:off-diagonal-empty-criterion}
For every polynomial map \(F\), the following are equivalent:
\[
  R_F^\circ=\varnothing
  \quad\Longleftrightarrow\quad
  C_F^\circ=0
  \quad\Longleftrightarrow\quad
  \Obs(F)=0
  \quad\Longleftrightarrow\quad
  I_R=I_\Delta.
\]
\end{corollary}

\begin{proof}
\[
  R_F^\circ=\varnothing
  \quad\Longleftrightarrow\quad
  C_F^\circ=0
  \quad\Longleftrightarrow\quad
  I_R:I_\Delta=S.
\]
Moreover,
\[
  I_R:I_\Delta=S
  \quad\Longleftrightarrow\quad
  I_\Delta\subseteq I_R
  \quad\Longleftrightarrow\quad
  I_R=I_\Delta.
\]
The equivalence with \(\Obs(F)=0\) is
Proposition~\ref{prop:obstruction-kernel}.
\end{proof}

The idempotent mechanism behind a clopen diagonal is purely
commutative-algebraic.  Let \(C\) be a commutative ring, \(J\subseteq C\)
an ideal, \(\delta\in C\), and \(c\in C^\times\).  Suppose
\[
  \delta J=0,
  \qquad
  \delta-c\in J.
\]
Define
\[
  q:=1-c^{-1}\delta.
\]

\begin{theorem}[Abstract off-diagonal projector]
\label{thm:abstract-off-diagonal-projector}
The element \(q\) is idempotent and
\[
  J=Cq,
  \qquad
  \Ann_C(J)=C(1-q).
\]
In particular,
\[
  J=0\quad\Longleftrightarrow\quad q=0.
\]
\end{theorem}

\begin{proof}
The congruence \(\delta\equiv c\pmod J\) gives \(q\in J\).  For every
\(j\in J\),
\[
  qj=j-c^{-1}\delta j=j.
\]
Taking \(j=q\) proves \(q^2=q\); the same identity gives \(J=Cq\).
Finally,
\[
  aq=0
  \quad\Longleftrightarrow\quad
  a=a(1-q),
\]
so \(\Ann_C(J)=\Ann_C(Cq)=C(1-q)\).
\end{proof}

Whenever the theorem is applied with
\[
  C=C_F,
  \qquad
  J=\Obs(F),
\]
the identities \(J=C_Fq\) and \(C_F/J\cong A\) give
\[
  C_F
  \cong
  A
  \times
  C_F/C_F(1-q).
\]
Thus the first factor is the diagonal and the second is its
off-diagonal complement.

\subsection{Generic fibers and marked-root decomposition}
\label{sec:generic-galois-bridge}

Let \(F\) be dominant and generically finite, and let
\[
  A=\C[x_1,\dots,x_n],
  \qquad
  B=\C[F_1,\dots,F_n]\subseteq A.
\]
Let
\[
  K=\Frac(B),
  \qquad
  L=\Frac(A).
\]
Let
\[
  \theta\colon K\otimes_BA\longrightarrow L,
  \qquad
  \lambda\otimes a\longmapsto\lambda a.
\]

\begin{theorem}[Generic collision bridge]
\label{thm:generic-collision-base-change}
The map \(\theta\) is an isomorphism.  Consequently, there is a
\(K\)-algebra isomorphism
\[
  K\otimes_BC_F
  \cong
  L\otimes_KL,
\]
and the following square commutes:
\[
  \begin{array}{ccc}
    K\otimes_BC_F
      &\xrightarrow{\ 1\otimes\bar\mu_F\ }&
    K\otimes_BA\\
    {\scriptstyle\sim}\downarrow
      &&\downarrow{\scriptstyle\theta\ \sim}\\
    L\otimes_KL
      &\xrightarrow{\ \mu_L\ }&
    L.
  \end{array}
\]
\end{theorem}

\begin{proof}
Let \(U=B\setminus\{0\}\).  Since \(K=U^{-1}B\), localization gives
\[
  K\otimes_BA
  \xrightarrow{\sim}
  U^{-1}A,
  \qquad
  \frac{b}{u}\otimes a\longmapsto\frac{ba}{u}.
\]
Because \(B\subseteq A\) and \(A\) is a domain,
\[
  U=B\setminus\{0\}\subseteq A\setminus\{0\},
  \qquad
  U^{-1}A\hookrightarrow\Frac(A)=L,
  \qquad
  \frac{a}{u}\longmapsto\frac{a}{u}.
\]
Composing this injection with the displayed isomorphism gives
\(\theta\), so \(\theta\) is injective.

Since \(F\) is generically finite, \(L/K\) is a finite extension.  The
image of \(U^{-1}A\) in \(L\) is
\[
  K[x_1,\dots,x_n].
\]
Each \(x_i\in L\) is algebraic over \(K\), so this is a
finite-dimensional \(K\)-algebra.  As a finite-dimensional
\(K\)-domain, it is a field: multiplication by a nonzero element is an
injective endomorphism of a finite-dimensional \(K\)-vector space and
therefore is surjective.
Because it contains \(A\), it contains \(\Frac(A)=L\); the reverse
inclusion is already given by the embedding above.  Therefore
\[
  U^{-1}A=L,
\]
and \(\theta\) is an isomorphism.

By Theorem~\ref{thm:collision-tensor-product}, identify
\(C_F\) with \(A\otimes_BA\).  There is a \(K\)-algebra isomorphism
\[
  \beta\colon
  K\otimes_B(A\otimes_BA)
  \xrightarrow{\sim}
  (K\otimes_BA)\otimes_K(K\otimes_BA)
\]
defined on pure tensors by
\[
  \beta\bigl(\lambda\otimes(a\otimes a')\bigr)
  =
  (\lambda\otimes a)\otimes(1\otimes a').
\]
Its inverse is
\[
  (\lambda\otimes a)\otimes(\lambda'\otimes a')
  \longmapsto
  \lambda\lambda'\otimes(a\otimes a').
\]
The balancing relations over \(B\) and \(K\) show that these formulas
are well-defined, and the two displayed maps are mutually inverse.
Composing \(\beta\) with \(\theta\otimes_K\theta\) gives
\[
  \Phi_F\colon K\otimes_BC_F\xrightarrow{\sim}L\otimes_KL.
\]

It remains to identify the diagonal map.  Under
\(C_F\cong A\otimes_BA\), the homomorphism \(\bar\mu_F\) is
multiplication.  Therefore, on a pure tensor,
\[
  \begin{aligned}
  (\mu_L\circ\Phi_F)
    \bigl(\lambda\otimes(a\otimes a')\bigr)
    &=(\lambda a)a'
     =\lambda aa',\\
  \bigl(\theta\circ(1\otimes\bar\mu_F)\bigr)
    \bigl(\lambda\otimes(a\otimes a')\bigr)
    &=\theta(\lambda\otimes aa')
     =\lambda aa'.
  \end{aligned}
\]
Pure tensors generate \(K\otimes_BC_F\), so the two homomorphisms are
equal.  This proves the commutative square.
\end{proof}

Suppose now that \(L/K\) is separable and has a power basis
\(L=K(\alpha)\).  Let \(m_\alpha(T)\in K[T]\) be the minimal polynomial
of \(\alpha\), and let \(h_\alpha(T)\in L[T]\) be determined by
\[
  m_\alpha(T)=(T-\alpha)h_\alpha(T)
  \qquad\text{in }L[T].
\]

\begin{theorem}[Marked-root tensor decomposition]
\label{thm:marked-root-tensor-decomposition}
Give \(L\otimes_KL\) its \(L\)-algebra structure through the left tensor
factor.  There is an \(L\)-algebra equivalence
\[
  L\otimes_KL
  \cong
  L\times L[T]/(h_\alpha),
\]
under which tensor multiplication \(L\otimes_KL\to L\) is projection
onto the first factor.
\end{theorem}

\begin{proof}
The power-basis presentation
\[
  L\cong K[T]/(m_\alpha)
\]
and base change from \(K\) to \(L\) give an \(L\)-algebra isomorphism
\[
  L\otimes_KL
  \xrightarrow{\sim}
  L[T]/(m_\alpha),
  \qquad
  1\otimes\alpha\longmapsto T.
\]
Since \(m_\alpha\) is separable,
\(m_\alpha'(\alpha)\neq0\).  Differentiating
\[
  m_\alpha(T)=(T-\alpha)h_\alpha(T)
\]
and evaluating at \(T=\alpha\) gives
\[
  h_\alpha(\alpha)=m_\alpha'(\alpha)\neq0.
\]
Thus \(T-\alpha\) does not divide \(h_\alpha\), so these two polynomials
are coprime in \(L[T]\).  The Chinese remainder theorem now gives
\[
  L[T]/(m_\alpha)
  \xrightarrow{\sim}
  L[T]/(T-\alpha)\times L[T]/(h_\alpha)
  \xrightarrow{\sim}
  L\times L[T]/(h_\alpha).
\]
Under the first isomorphism above, a pure tensor
\(\ell\otimes p(\alpha)\), with \(p\in K[T]\), corresponds to the class
of \(\ell p(T)\).  Tensor multiplication sends it to
\(\ell p(\alpha)\), which is evaluation at \(T=\alpha\).  Under the
Chinese remainder isomorphism, this evaluation map is projection onto
the factor \(L[T]/(T-\alpha)\cong L\).  Hence tensor multiplication is
the first projection.
\end{proof}

Together with Theorem~\ref{thm:generic-collision-base-change}, this gives
\[
  K\otimes_BC_F
  \cong
  L\times L[T]/(h_\alpha),
\]
where the first projection is the generic diagonal.

\subsection{Normalized covers and inertia}
\label{sec:normalized-covers-inertia}

Let \(L/K\) be finite and separable, and let \(N/K\) be its normal
closure, with a fixed embedding \(L\hookrightarrow N\).  Let
\[
  G=\Gal(N/K),
  \qquad
  H=\Gal(N/L).
\]

\begin{proposition}[Conjugate-embedding action]
\label{prop:conjugate-embedding-action}
Restriction induces a bijection
\[
  G/H\xrightarrow{\sim}\Hom_K(L,N),
  \qquad
  gH\longmapsto g|_L.
\]
The induced action of \(G\) on \(G/H\) is faithful; equivalently,
\[
  \core_G(H)=\bigcap_{g\in G}gHg^{-1}=1.
\]
\end{proposition}

\begin{proof}
For \(g_1,g_2\in G\),
\[
  g_1|_L=g_2|_L
  \quad\Longleftrightarrow\quad
  g_2^{-1}g_1\in H
  \quad\Longleftrightarrow\quad
  g_1H=g_2H.
\]
Every \(K\)-embedding \(L\hookrightarrow N\) extends to an element of
\(G\), so restriction gives the stated bijection.  Its action kernel is
\[
  \ker\!\left(G\curvearrowright G/H\right)
  =
  \bigcap_{g\in G}gHg^{-1}
  =
  \core_G(H).
\]
If \(\sigma\in\core_G(H)\), then
\[
  \sigma|_{g(L)}=\id_{g(L)}
  \quad(g\in G).
\]
Since \(N\) is the normal closure,
\[
  N=K\bigl(g(L):g\in G\bigr),
\]
and therefore \(\sigma=\id_N\).
\end{proof}

Let
\[
  Y=\Spec B,
  \qquad
  \overline X=\Norm_L(Y),
  \qquad
  Z=\Norm_N(Y).
\]
The marked embedding \(L\hookrightarrow N\) induces a canonical
commutative triangle
\[
  Z\xrightarrow{\pi}\overline X\xrightarrow{\bar\nu}Y,
  \qquad
  \nu=\bar\nu\circ\pi\colon Z\longrightarrow Y.
\]
Suppose that \(\bar\nu\) and \(\nu\) are finite, and fix an open
immersion over \(Y\)
\[
  j\colon X=\Spec A\hookrightarrow\overline X
\]
with \(\bar\nu\circ j=f_F\).  Denote this diagram and its deleted boundary
by
\[
  \mathcal D=
  \left(
    Z\xrightarrow{\pi}\overline X\xrightarrow{\bar\nu}Y,\,
    j\colon X\hookrightarrow\overline X
  \right),
  \qquad
  \partial X=\overline X\setminus j(X).
\]

Fix \(E\in Z^{(1)}\).  Let \(w_E\) be its divisorial
valuation, let \(b=\nu(E)\), and let \(v_b=w_E|_K\).
Let \(D_E\leq G\) and \(I_E\trianglelefteq D_E\) be its decomposition and
inertia groups.  Let
\[
  \mathcal Q_E:=D_E\backslash G/H.
\]
For \(\omega=D_EgH\in\mathcal Q_E\), let
\[
  u_\omega=(g^{-1}w_E)|_L,
  \qquad
  \operatorname{cent}(E,\omega)
  =\operatorname{cent}_{\overline X}(u_\omega),
\]
and let
\[
  i_E(\omega)=[I_E:I_E\cap gHg^{-1}].
\]
\begin{proposition}[Conjugate divisors and double cosets]
\label{prop:conjugate-divisor-double-cosets}
The assignments above are well-defined, and restriction of valuations
induces a bijection
\[
  \mathcal Q_E
  \xrightarrow{\sim}
  \{\text{prime divisors of }\overline X\text{ above }b\},
  \qquad
  \omega\longmapsto\operatorname{cent}(E,\omega).
\]
\end{proposition}

\begin{proof}
The group \(G\) acts transitively on the extensions of \(v_b\) to \(N\),
and the stabilizer of \(w_E\) is \(D_E\).  Restricting such a valuation
to \(L=N^H\) identifies precisely those extensions whose indices differ
on the right by an element of \(H\).  The resulting orbits are therefore
indexed by \(D_E\backslash G/H\), giving the stated bijection.

It remains to check independence of representatives.  If
\(D_Eg'H=D_EgH\), then \(g'=dgh\) for some \(d\in D_E\) and \(h\in H\).
The element \(d\) fixes \(w_E\), while \(h\) fixes \(L\) pointwise, so
\[
  (g'^{-1}w_E)|_L=(g^{-1}w_E)|_L.
\]
Moreover, \(I_E\trianglelefteq D_E\), and conjugation by \(d\) identifies
\[
  I_E\cap g'Hg'^{-1}
  \quad\text{with}\quad
  I_E\cap gHg^{-1}.
\]
Thus \(u_\omega\), \(\operatorname{cent}(E,\omega)\), and
\(i_E(\omega)\) are
well-defined.
\end{proof}

We use the notation
\[
  e\bigl(\operatorname{cent}(E,\omega)/b\bigr):=e(u_\omega/v_b).
\]
Since the residue characteristic is zero,
\[
  \Ram^{(1)}(Z/Y)
  :=
  \{E\in Z^{(1)}:e(w_E/v_b)>1\}
  =
  \{E\in Z^{(1)}:I_E\neq1\}.
\]

\begin{proposition}[Detection of inertia on a conjugate sheet]
\label{prop:core-free-sheet-action}
If \(I_E\neq1\), then there is an \(\omega\in\mathcal Q_E\) with
\[
  i_E(\omega)>1.
\]
\end{proposition}

\begin{proof}
Suppose, to the contrary, that \(i_E(\omega)=1\) for every
\(\omega\in\mathcal Q_E\).  For each \(g\in G\),
\[
  1=i_E(D_EgH)=[I_E:I_E\cap gHg^{-1}],
\]
so
\[
  I_E=I_E\cap gHg^{-1}\subseteq gHg^{-1}
  \qquad\text{for every }g\in G.
\]
Consequently,
\[
  I_E
  \subseteq
  \bigcap_{g\in G}gHg^{-1}
  =
  \core_G(H)
  =
  1,
\]
contrary to \(I_E\neq1\).
\end{proof}

\begin{proposition}[Intermediate ramification index]
\label{prop:intermediate-ramification-index}
For \(\omega=D_EgH\in\mathcal Q_E\),
\[
  e\bigl(\operatorname{cent}(E,\omega)/b\bigr)
  =
  i_E(\omega)
  =
  [I_E:I_E\cap gHg^{-1}].
\]
\end{proposition}

\begin{proof}
Let \(w_E\) and \(v_b\) be the divisorial valuations associated with
\(E\) and \(b\), and let
\[
  w_g=g^{-1}w_E,
  \qquad
  u_\omega=w_g|_L.
\]
The valuation \(u_\omega\) has center \(\operatorname{cent}(E,\omega)\) on
\(\overline X\).  Since
the residue fields have characteristic zero,
\[
  e(w_g/v_b)=|I_E|.
\]
The inertia group of \(w_g\) over \(u_\omega\) is
\[
  g^{-1}I_Eg\cap H,
\]
and hence
\[
  e(w_g/u_\omega)
  =
  |g^{-1}I_Eg\cap H|
  =
  |I_E\cap gHg^{-1}|.
\]
Multiplicativity of ramification indices in the field tower
\[
  K\subseteq L\subseteq N
\]
therefore gives
\[
  e(u_\omega/v_b)
  =
  \frac{e(w_g/v_b)}{e(w_g/u_\omega)}
  =
  [I_E:I_E\cap gHg^{-1}],
\]
which is the asserted formula.
\end{proof}

\begin{proposition}[Visibility-to-boundary]
\label{prop:visibility-to-boundary}
Suppose that \(X\to Y\) is \'etale.  For \(E\in Z^{(1)}\) and
\(\omega\in\mathcal Q_E\),
\[
  i_E(\omega)>1
  \quad\Longrightarrow\quad
  \operatorname{cent}(E,\omega)\in\partial X.
\]
\end{proposition}

\begin{proof}
Suppose that \(\operatorname{cent}(E,\omega)\in j(X)\), and let
\(x\in X\) be the corresponding point under the open immersion \(j\).
The local homomorphism
\[
  \mathcal O_{Y,b}
  \longrightarrow
  \mathcal O_{X,x}
\]
is \'etale, so
\(e(\operatorname{cent}(E,\omega)/b)=1\).  Proposition
\ref{prop:intermediate-ramification-index} now gives
\[
  i_E(\omega)
  =
  e\bigl(\operatorname{cent}(E,\omega)/b\bigr)
  =
  1.
\]
Taking the contrapositive gives the result.
\end{proof}

\begin{definition}[Hidden-inertia locus]
\label{def:no-hidden-inertia}
\[
  \mathcal H(\mathcal D)
  :=
  \left\{
  E\in\Ram^{(1)}(Z/Y):
  \begin{array}{l}
    \exists\omega\in\mathcal Q_E,\ i_E(\omega)>1,\\
    \forall\omega\in\mathcal Q_E,\quad
      i_E(\omega)>1
      \Longrightarrow
      \operatorname{cent}(E,\omega)\in\partial X
  \end{array}
  \right\}.
\]
Thus \(E\in\mathcal H(\mathcal D)\) precisely when nontrivial inertia
is visible on a conjugate sheet but every visible center belongs to the
deleted boundary.
\end{definition}

\subsection{Automorphism detection}
\label{sec:automorphism-detection}

\begin{theorem}[Automorphism criterion]
\label{thm:automorphism-criterion}
For an arbitrary polynomial self-map
\(F\colon\A^n_{\C}\to\A^n_{\C}\), the following are equivalent:
\begin{enumerate}[label=\textup{(\arabic*)}]
  \item \(F\) is a polynomial automorphism;
  \item \(F\) is injective on \(\A^n(\C)\);
  \item \(I_R(F)=I_\Delta\);
  \item \(\Obs(F)=0\).
\end{enumerate}
\end{theorem}

\begin{proof}
Suppose first that \(F\) has polynomial inverse
\(G=(G_1,\dots,G_n)\).  Then
\[
  x_i-y_i
  =
  G_i(F(x))-G_i(F(y))
  \in I_R
\]
for every \(i\), so \(I_\Delta\subseteq I_R\).  Proposition
\ref{prop:collision-contained-in-diagonal} gives equality.  Thus
\textup{(1)} implies \textup{(3)}.

Suppose \(I_R=I_\Delta\), and let
\(\operatorname{ev}_{(a,b)}\colon S\to\C\) be evaluation at
\((a,b)\).  If \(F(a)=F(b)\), then
\[
  I_R\subseteq\ker(\operatorname{ev}_{(a,b)}).
\]
Therefore \(x_i-y_i\in I_\Delta=I_R\) gives \(a_i-b_i=0\) for every
\(i\).  Thus \(a=b\), proving
\textup{(3)}\(\Rightarrow\)\textup{(2)}.  The
Ax--Grothendieck theorem \cite{Ax1969} gives
\textup{(2)}\(\Rightarrow\)\textup{(1)}.  Finally,
Proposition~\ref{prop:obstruction-kernel} gives
\textup{(3)}\(\Longleftrightarrow\)\textup{(4)}.
\end{proof}

\begin{corollary}[The Jacobian conjecture as obstruction vanishing]
\label{cor:jacobian-conjecture-as-vanishing}
For every \(n\geq1\), the \(n\)-dimensional Jacobian conjecture is
equivalent to each of the statements
\[
  \forall F\colon\A^n_{\C}\to\A^n_{\C},
  \qquad
  \det JF\in\C^\times
  \quad\Longrightarrow\quad
  \Obs(F)=0,
\]
and
\[
  \forall F\colon\A^n_{\C}\to\A^n_{\C},
  \qquad
  \det JF\in\C^\times
  \quad\Longrightarrow\quad
  I_R(F)=I_\Delta.
\]
\end{corollary}

\begin{proof}
Apply Theorem~\ref{thm:automorphism-criterion} to every Keller map.
Its equivalent conditions \textup{(1)}, \textup{(3)}, and \textup{(4)}
give the two displayed formulations.
\end{proof}

\subsection{Generic degree one}

\begin{corollary}[Generic degree-one descent]
\label{cor:generic-degree-one-descent}
If
\[
  L=K
  \qquad\text{and}\qquad
  A\text{ is flat over }B,
\]
then
\[
  \Obs(F)=0.
\]
\end{corollary}

\begin{proof}
When \(L=K\), Theorem~\ref{thm:generic-collision-base-change}
identifies the base-changed diagonal with tensor multiplication:
\[
  \ker(1\otimes\bar\mu_F)
  \cong
  \ker\!\left(
    K\otimes_KK\xrightarrow{\mu_K}K
  \right)
  =
  0.
\]
The fiber-product identification of
Theorem~\ref{thm:collision-tensor-product} is
\[
  C_F\xrightarrow{\sim}A\otimes_BA.
\]
Since \(A\) is flat over \(B\), the tensor product \(A\otimes_BA\), and
therefore \(C_F\), is flat over \(B\).  The localization map
\[
  \iota_C\colon C_F\hookrightarrow K\otimes_BC_F
\]
is therefore injective, and
\[
  \iota_C\bigl(\ker\bar\mu_F\bigr)
  \subseteq
  \ker(1\otimes\bar\mu_F)
  =
  0.
\]
Thus \(\ker\bar\mu_F=0\).  Proposition
\ref{prop:obstruction-kernel} identifies
\[
  \ker(\bar\mu_F)=\Obs(F)=I_\Delta/I_R,
\]
so \(\Obs(F)=0\).
\end{proof}

\begin{leanfiles}
\leanpath{CollisionIdeals/CollisionDiagonal.lean},
\leanpath{CollisionIdeals/OffDiagonalScheme.lean},
\leanpath{CollisionIdeals/Secant.lean},
\leanpath{CollisionIdeals/GenericCollisionDiagonal.lean},
\leanpath{CollisionIdeals/KellerCollisionModel.lean}, and
\leanpath{CollisionIdeals/JacobianConjecture.lean}.  The Lean form of
Theorem~\ref{thm:generic-collision-base-change} takes surjectivity of
\(\theta\) as a parameter; the proof above derives it from generic
finiteness.
\end{leanfiles}

\section{Planar vanishing}
\label{sec:planar-vanishing}

Let
\[
  F=(P,Q)\colon\A^2_{\C}\longrightarrow\A^2_{\C}
\]
be a polynomial map satisfying
\(\det JF=c\in\C^\times\).

\subsection{The planar secant projector}

Let \(\Delta=V(I_\Delta)\subseteq\Spec S\).  For \(a\in S\), write
\(a|_\Delta\) for its image under \(S\to S/I_\Delta\), and apply this
notation entrywise to matrices.
Let
\[
  \mathbf d=
  \begin{pmatrix}
    x_1-y_1\\
    x_2-y_2
  \end{pmatrix},
  \qquad
  \mathbf r=
  \begin{pmatrix}
    P(x)-P(y)\\
    Q(x)-Q(y)
  \end{pmatrix}.
\]
For \(H\in\C[x_1,x_2]\), let
\[
  D_1H
  =
  \frac{H(x_1,x_2)-H(y_1,x_2)}{x_1-y_1},
  \qquad
  D_2H
  =
  \frac{H(y_1,x_2)-H(y_1,y_2)}{x_2-y_2}.
\]
Both quotients belong to \(S\).  Let
\[
  M_F=
  \begin{pmatrix}
    D_1P&D_2P\\
    D_1Q&D_2Q
  \end{pmatrix}.
\]
Then
\[
  \mathbf r=M_F\mathbf d,
  \qquad
  M_F\big|_\Delta=JF.
\]
Let \(\delta_F=\det(M_F)\).

\begin{proposition}[Planar secant annihilation]
\label{prop:secant-annihilator}
The secant determinant satisfies
\[
  \delta_F I_\Delta\subseteq I_R,
  \qquad
  \delta_F\big|_\Delta=c.
\]
\end{proposition}

\begin{proof}
Multiplying \(\mathbf r=M_F\mathbf d\) by
\(\operatorname{adj}(M_F)\) gives
\[
  \delta_F\mathbf d=\operatorname{adj}(M_F)\mathbf r.
\]
\[
  \mathbf r\in I_R^{\oplus2},
  \qquad
  \operatorname{adj}(M_F)\in M_2(S)
  \quad\Longrightarrow\quad
  \operatorname{adj}(M_F)\mathbf r\in I_R^{\oplus2}.
\]
Therefore
\[
  \delta_F\mathbf d\in I_R^{\oplus2},
\]
or equivalently,
\[
  \forall j\in\{1,2\},
  \qquad
  \delta_F(x_j-y_j)\in I_R.
\]
Hence \(\delta_FI_\Delta\subseteq I_R\).  On the diagonal,
\[
  (D_jF_i)|_\Delta=\frac{\partial F_i}{\partial x_j},
\]
so \(M_F|_\Delta=JF\) and
\(\delta_F|_\Delta=\det JF=c\).
\end{proof}

Let \(\bar\delta_F\) be the class of \(\delta_F\) in \(C_F\), and define
\[
  q_F:=1-c^{-1}\bar\delta_F.
\]

\begin{theorem}[Planar collision-sheet decomposition]
\label{thm:collision-sheet-decomposition}
The collision scheme decomposes as
\[
  R_F\cong\Delta_{X/Y}(X)\amalg R_F^\circ,
  \qquad
  C_F\cong A\times C_F^\circ.
\]
The off-diagonal ideal satisfies
\[
  I_{\mathrm{off}}
  =
  I_R+(\delta_F)
  =
  I_R:I_\Delta^\infty.
\]
\end{theorem}

\begin{proof}
By Proposition~\ref{prop:obstruction-kernel},
\[
  J
  :=
  \ker(\bar\mu_F)
  =
  I_\Delta/I_R
  =
  \Obs(F).
\]
Proposition~\ref{prop:secant-annihilator} gives
\[
  \bar\delta_FJ=0,
  \qquad
  \bar\delta_F-c\in J.
\]
Theorem~\ref{thm:abstract-off-diagonal-projector}, applied to these two
relations and \(q_F=1-c^{-1}\bar\delta_F\), yields
\[
  q_F^2=q_F,
  \qquad
  J=C_Fq_F,
  \qquad
  \Ann_{C_F}(J)=C_F(1-q_F),
  \qquad
  J=0\Longleftrightarrow q_F=0.
\]
Because \(1-q_F=c^{-1}\bar\delta_F\),
\[
  \frac{I_R:I_\Delta}{I_R}
  =
  \Ann_{C_F}(J)
  =
  C_F(1-q_F)
  =
  C_F\bar\delta_F
  =
  \frac{I_R+(\delta_F)}{I_R}.
\]
Hence
\[
  I_{\mathrm{off}}
  =
  I_R:I_\Delta
  =
  I_R+(\delta_F).
\]

Idempotence gives \(J^m=C_Fq_F=J\) for every \(m\geq1\).  Therefore
\[
  \frac{I_R:I_\Delta^m}{I_R}
  =
  \Ann_{C_F}(J^m)
  =
  \Ann_{C_F}(J)
  =
  \frac{I_R:I_\Delta}{I_R}.
\]
Thus
\[
  I_R:I_\Delta^\infty
  =
  \bigcup_{m\geq1}(I_R:I_\Delta^m)
  =
  I_R:I_\Delta.
\]

The idempotent \(q_F\) also gives a ring isomorphism
\[
  C_F
  \xrightarrow{\ \sim\ }
  C_F(1-q_F)\times C_Fq_F,
  \qquad
  a\longmapsto\bigl(a(1-q_F),aq_F\bigr),
\]
whose inverse is addition.  The two factors satisfy
\[
  C_F(1-q_F)
  \cong
  C_F/C_Fq_F
  =
  C_F/J
  \xrightarrow[\bar\mu_F]{\ \sim\ }
  A
\]
and
\[
  C_Fq_F
  \cong
  C_F/C_F(1-q_F)
  \cong
  S\big/\bigl(I_R+(\delta_F)\bigr)
  =
  C_F^\circ.
\]
Thus \(C_F\cong A\times C_F^\circ\).  Since the spectrum of a product is
the disjoint union of the two spectra, this isomorphism dualizes to
\[
  R_F\cong\Delta_{X/Y}(X)\amalg R_F^\circ.
\]
\end{proof}

Corollary~\ref{cor:off-diagonal-empty-criterion} states that
\[
  R_F^\circ=\varnothing
  \quad\Longleftrightarrow\quad
  C_F^\circ=0
  \quad\Longleftrightarrow\quad
  \Obs(F)=0
  \quad\Longleftrightarrow\quad
  I_R=I_\Delta.
\]
The preceding theorem identifies
\[
  C_F^\circ
  =
  S\big/\bigl(I_R+(\delta_F)\bigr),
\]
and Theorem~\ref{thm:abstract-off-diagonal-projector} gives
\(\Obs(F)=0\Longleftrightarrow q_F=0\).  Hence
\[
  q_F=0
  \quad\Longleftrightarrow\quad
  I_R+(\delta_F)=S
  \quad\Longleftrightarrow\quad
  R_F^\circ=\varnothing
  \quad\Longleftrightarrow\quad
  I_R=I_\Delta.
\]
In particular, the secant determinant cuts out the off-diagonal factor:
\[
  R_F^\circ
  =
  \Spec\!\bigl(S/(I_R+(\delta_F))\bigr).
\]

\subsection{Boundary exclusion and inertia}

Let
\[
  A=\C[x_1,x_2],
  \qquad
  B=\C[P,Q]\subseteq A,
  K=\C(P,Q),
  \qquad
  L=\C(x_1,x_2),
\]
and let \(X=\Spec A\) and \(Y=\Spec B\).

\begin{proposition}[Keller local geometry]
\label{prop:planar-keller-local-geometry}
\[
  \C[u,v]\xrightarrow{\sim}B,
  \qquad
  u\longmapsto P,\quad v\longmapsto Q.
\]
Moreover,
\[
  X\xrightarrow{\ f_F\ }Y\ \text{is \'etale},
  \qquad
  A\ \text{is flat over }B,
  \qquad
  [L:K]<\infty
  \ \text{and }L/K\text{ is separable}.
\]
\end{proposition}

\begin{proof}
The identity
\[
  dP\wedge dQ
  =
  \left(
    \frac{\partial P}{\partial x_1}
    \frac{\partial Q}{\partial x_2}
    -
    \frac{\partial P}{\partial x_2}
    \frac{\partial Q}{\partial x_1}
  \right)dx_1\wedge dx_2
  =
  \det(JF)\,dx_1\wedge dx_2
  =
  c\,dx_1\wedge dx_2
  \neq0
\]
shows that \(P,Q\) are algebraically independent.  Thus
\(B\cong\C[u,v]\) and \(Y\cong\A^2_{\C}\).  Moreover,
\[
  A
  \cong
  B[U,V]\big/\bigl(P(U,V)-u,\ Q(U,V)-v\bigr),
\]
whose relative Jacobian determinant is \(c\in A^\times\).  The
Jacobian criterion therefore makes \(B\to A\) \'etale.  In particular
it is flat and quasi-finite.  Since \(f_F\) is dominant and
quasi-finite, its generic fiber is zero-dimensional, so
\([L:K]<\infty\).  The morphism remains \'etale at the generic point;
hence the finite extension \(L/K\) is separable.
\end{proof}

Fix an algebraic closure \(\overline K\supseteq L\), and let
\[
  \begin{aligned}
  N
    &:=
      K\bigl(\sigma(L):
        \sigma\in\Hom_K(L,\overline K)\bigr),
    &K&\subseteq L\subseteq N\subseteq\overline K,\\
  \overline A
    &:=
      \{a\in L:a\text{ is integral over }B\},
  &
  A_N
    &:=
      \{a\in N:a\text{ is integral over }B\},\\
  \overline X&:=\Spec\overline A,
  &
  Z&:=\Spec A_N.
  \end{aligned}
\]
Thus \(N/K\) is the normal closure of \(L/K\).

\begin{proposition}[Planar finite-cover diagram]
\label{prop:planar-normalization-diagram}
\[
  Z\xrightarrow[\mathrm{finite}]{\pi}
  \overline X\xrightarrow[\mathrm{finite}]{\bar\nu}Y,
  \qquad
  \nu=\bar\nu\circ\pi,
\]
and
\[
  X\xhookrightarrow[\mathrm{open}]{j}\overline X,
  \qquad
  \bar\nu\circ j=f_F.
\]
\end{proposition}

\begin{proof}
Since \(B\cong\C[u,v]\) is a finitely generated normal
\(\C\)-algebra, its integral closures in the finite extensions \(L/K\) and
\(N/K\) are finite over \(B\).  The inclusion \(L\subseteq N\) gives
\(\overline A\subseteq A_N\), so
\[
  B\subseteq\overline A\subseteq A_N.
\]
The morphism \(\bar\nu\) is finite because \(\overline A\) is a finite
\(B\)-module.  If \(a_1,\dots,a_r\) generate \(A_N\) over \(B\), then
they also generate \(A_N\) over \(\overline A\), since
\(B\subseteq\overline A\).  Hence \(\pi\) is finite.

For \(a\in\overline A\), integrality over \(B\subseteq A\) and the
normality of \(A=\C[x_1,x_2]\) give \(a\in A\).  Thus
\[
  B\subseteq\overline A\subseteq A,
  \qquad
  \Frac(\overline A)=\Frac(A)=L,
\]
induces a dominant birational morphism
\(j\colon X\to\overline X\).  It is quasi-finite and separated by
Proposition~\ref{prop:planar-keller-local-geometry}.  Since
\(\overline X\) is normal, Zariski's Main Theorem makes \(j\) an open
immersion.
\end{proof}

Let
\[
  \begin{gathered}
    \mathcal D_F=
    \left(
      Z\xrightarrow{\pi}\overline X\xrightarrow{\bar\nu}Y,\,
      j\colon X\hookrightarrow\overline X
    \right),
    \qquad
    \partial X=\overline X\setminus j(X),\\
    G=\Gal(N/K),
    \qquad
    H=\Gal(N/L).
  \end{gathered}
\]

For \(E\in\Ram^{(1)}(Z/Y)\) and
\(\omega\in\mathcal Q_E\), Proposition
\ref{prop:intermediate-ramification-index} gives
\[
  e\bigl(\operatorname{cent}(E,\omega)/\nu(E)\bigr)=i_E(\omega).
\]
Since \(X\to Y\) is \'etale by
Proposition~\ref{prop:planar-keller-local-geometry},
Proposition~\ref{prop:visibility-to-boundary} gives
\[
  i_E(\omega)>1
  \quad\Longrightarrow\quad
  \operatorname{cent}(E,\omega)\in\partial X.
\]
Thus Definition~\ref{def:no-hidden-inertia} specializes to
\[
  \mathcal H(\mathcal D_F)
  =
  \left\{
  E\in\Ram^{(1)}(Z/Y):
  \begin{array}{l}
    \exists\omega\in\mathcal Q_E,\ i_E(\omega)>1,\\
    \forall\omega\in\mathcal Q_E,\quad
      i_E(\omega)>1
      \Longrightarrow
      \operatorname{cent}(E,\omega)\in\partial X
  \end{array}
  \right\}.
\]

\subsection{Hidden inertia and divisorial ramification}
\label{sec:planar-no-hidden-inertia}

\begin{corollary}[Planar divisorial ramification criterion]
\label{cor:planar-codimension-one-unramified}
\[
  \mathcal H(\mathcal D_F)=\varnothing
  \quad\Longleftrightarrow\quad
  \Ram^{(1)}(Z/Y)=\varnothing.
\]
\end{corollary}

\begin{proof}
The reverse implication follows from
\(\mathcal H(\mathcal D_F)\subseteq\Ram^{(1)}(Z/Y)\).  Conversely,
suppose that \(\mathcal H(\mathcal D_F)=\varnothing\) and let
\(E\in\Ram^{(1)}(Z/Y)\).  Then \(I_E\neq1\).
Since \(N\) is the normal closure of \(L/K\),
\(\core_G(H)=1\), and
Proposition~\ref{prop:core-free-sheet-action} gives
\[
  \exists\omega\in\mathcal Q_E,
  \qquad
  i_E(\omega)>1.
\]
For every \(\omega'\in\mathcal Q_E\), Proposition
\ref{prop:visibility-to-boundary} gives
\[
  i_E(\omega')>1
  \quad\Longrightarrow\quad
  \operatorname{cent}(E,\omega')\in\partial X.
\]
Hence \(E\in\mathcal H(\mathcal D_F)\), a contradiction.
\end{proof}

\subsection{Collapse of the normal closure}

\begin{theorem}[Purity of the branch locus, specialized]
\label{thm:branch-purity}
Let \(Z\to\A^2_{\C}\) be a finite dominant generically separable
morphism with \(Z\) integral and normal.  If
\(\Ram^{(1)}(Z/\A^2_{\C})=\varnothing\), then the morphism is \'etale
\cite{StacksPurity}.
\end{theorem}

\begin{theorem}[Finite-\'etale rigidity of the affine plane]
\label{thm:a2-etale-rigidity}
Every connected finite \'etale cover of \(\A^2_{\C}\) is trivial
\cite[Expos\'e XII]{SGA1}.
\end{theorem}

\begin{corollary}[Triviality of the planar normal closure]
\label{cor:planar-normal-closure-trivial}
If
\[
  \Ram^{(1)}(Z/Y)=\varnothing,
\]
then
\[
  N=K
  \qquad\text{and hence}\qquad
  L=K.
\]
\end{corollary}

\begin{proof}
By Proposition~\ref{prop:planar-keller-local-geometry},
\(Y\cong\A^2_{\C}\).  Since \(N/K\) is finite Galois,
\(\nu\colon Z\to Y\) is generically separable; hence
Theorem~\ref{thm:branch-purity} makes it finite \'etale.  The inclusion
\[
  A_N\subset N\text{ is a domain},
  \qquad\text{so}\qquad
  Z=\Spec A_N\text{ is connected}.
\]
Theorem~\ref{thm:a2-etale-rigidity} therefore makes
\(\nu\colon Z\to Y\) an isomorphism.  Taking function fields gives
\(N=K\).  Finally,
\(K\subseteq L\subseteq N\) implies \(L=K\).
\end{proof}

\subsection{Vanishing of the off-diagonal factor}

\begin{theorem}[Planar Vanishing]
\label{thm:planar-vanishing}
For every polynomial map
\(F=(P,Q)\colon\A^2_{\C}\to\A^2_{\C}\) satisfying
\(\det JF=c\in\C^\times\),
\[
  \mathcal H(\mathcal D_F)=\varnothing
  \quad\Longrightarrow\quad
  \begin{gathered}
    I_R=I_\Delta,
    \qquad
    \Obs(F)=0,\\
    q_F=0,
    \qquad
    R_F^\circ=\varnothing.
  \end{gathered}
\]
\end{theorem}

\begin{proof}
Corollary~\ref{cor:planar-codimension-one-unramified} gives
\(\Ram^{(1)}(Z/Y)=\varnothing\), and
Corollary~\ref{cor:planar-normal-closure-trivial} gives \(L=K\).
Proposition~\ref{prop:planar-keller-local-geometry} gives that \(A\) is
flat over \(B\).  Corollary~\ref{cor:generic-degree-one-descent} states
\[
  L=K,
  \qquad
  A\text{ flat over }B
  \quad\Longrightarrow\quad
  \Obs(F)=0,
\]
so \(\Obs(F)=0\).

Corollary~\ref{cor:off-diagonal-empty-criterion} and
Theorem~\ref{thm:abstract-off-diagonal-projector} state
\[
  \begin{gathered}
    R_F^\circ=\varnothing
    \quad\Longleftrightarrow\quad
    C_F^\circ=0
    \quad\Longleftrightarrow\quad
    \Obs(F)=0
    \quad\Longleftrightarrow\quad
    I_R=I_\Delta,\\
    \Obs(F)=0
    \quad\Longleftrightarrow\quad
    q_F=0.
  \end{gathered}
\]
All four conclusions follow.
\end{proof}

\begin{corollary}[Conditional Jacobian conjecture in dimension two]
\label{cor:conditional-jc2}
If \(\mathcal H(\mathcal D_F)=\varnothing\) for every planar Keller map
\(F\), then every planar Keller map is a polynomial automorphism.
Equivalently, the two-dimensional Jacobian conjecture holds.
\end{corollary}

\begin{proof}
The automorphism criterion, Theorem~\ref{thm:automorphism-criterion},
states that every polynomial self-map of \(\A^2_{\C}\) satisfies
\[
  I_R=I_\Delta
  \quad\Longleftrightarrow\quad
  F\text{ is a polynomial automorphism}.
\]
Planar Vanishing supplies the left-hand side for every planar Keller
map under the stated condition.
\end{proof}

\paragraph{External inputs to the formalization.}
The Lean development takes the following established results as
explicit inputs:
\begin{enumerate}[
  label=\textup{(E\arabic*)},
  leftmargin=*,
  itemsep=2pt,
  topsep=3pt
]
  \item The planar Ax--Grothendieck theorem: every injective polynomial
    self-map of \(\A^2_{\C}\) is a polynomial automorphism
    \cite{Ax1969}.
  \item Purity of the branch locus in the specialized form of
    Theorem~\ref{thm:branch-purity} \cite{StacksPurity}.
  \item Every connected finite \'etale cover of \(\A^2_{\C}\) is
    trivial, as in Theorem~\ref{thm:a2-etale-rigidity}
    \cite[Expos\'e XII]{SGA1}.
\end{enumerate}
The condition \(\mathcal H(\mathcal D_F)=\varnothing\) is not included
among these inputs; it remains the explicit premise of
Theorem~\ref{thm:planar-vanishing}.

\begin{leanfiles}
The conditional planar chain is collected in
\leanpath{CollisionIdeals/Planar.lean}.  The three inputs above are
isolated in
\leanpath{CollisionIdeals/Planar/ExternalAssumptions.lean}.
\end{leanfiles}

\section{\texorpdfstring{The cubic \(S_3\) obstruction in dimension three}
  {The cubic S3 obstruction in dimension three}}
\label{sec:cubic-s3-obstruction}

Let \(S_3=\operatorname{Perm}(\{1,2,3\})\) be the symmetric group on
three letters.

\begin{proposition}[Dimension-three counterexample criterion]
\label{prop:dimension-three-counterexample-criterion}
For a polynomial map
\[
  F\colon\A^3_{\C}\longrightarrow\A^3_{\C},
\]
consider the following statements:
\begin{enumerate}[label=\textup{(\arabic*)}]
  \item \(\det JF\in\C^\times\);
  \item \(F\) is not a polynomial automorphism;
  \item \(R_F^\circ\neq\varnothing\).
\end{enumerate}
Under \textup{(1)}, statements \textup{(2)} and \textup{(3)} are
equivalent.
\end{proposition}

\begin{proof}
By Corollary~\ref{cor:off-diagonal-empty-criterion} and
Theorem~\ref{thm:automorphism-criterion},
\[
  R_F^\circ=\varnothing
  \quad\Longleftrightarrow\quad
  I_R=I_\Delta
  \quad\Longleftrightarrow\quad
  F\in\Aut_{\mathrm{poly}}(\A^3_{\C}).
\]
Negating this equivalence proves
\textup{(2)}\(\Longleftrightarrow\)\textup{(3)}.
\end{proof}

\subsection{The cubic extension and its normal closure}
\label{sec:cubic-extension}

Let
\[
  F=(F_1,F_2,F_3)\colon\A^3_{\C}\longrightarrow\A^3_{\C}
\]
be dominant and generically finite.  Let
\[
  A=\C[x_1,x_2,x_3],
  \qquad
  B=\C[F_1,F_2,F_3]\subseteq A,
  \qquad
  K=\Frac(B),
  \qquad
  L=\Frac(A),
\]
and let
\[
  \theta\colon K\otimes_BA\longrightarrow L,
  \qquad
  \lambda\otimes a\longmapsto\lambda a,
\]
be the isomorphism of
Theorem~\ref{thm:generic-collision-base-change}.  Suppose that \(L/K\)
is separable and nonnormal, and
choose a power basis with generator \(\alpha\) and degree three:
\[
  L=K(\alpha),
  \qquad
  [L:K]=3.
\]
If \(m_\alpha\in K[T]\) is the minimal polynomial of \(\alpha\), let
\(h_\alpha\in L[T]\) be determined by
\[
  m_\alpha(T)=(T-\alpha)h_\alpha(T)
  \qquad\text{in }L[T].
\]
Let \(N/K\) be the normal closure of \(L/K\), containing \(L\) through
the marked inclusion.  Let
\[
  G=\Gal(N/K),
  \qquad
  H=\Gal(N/L).
\]
Recall from Section~\ref{sec:normalized-covers-inertia} that
\[
  G/H\xrightarrow{\sim}\Hom_K(L,N),
  \qquad
  gH\longmapsto g|_L,
  \qquad
  \core_G(H)=1.
\]
The final equality follows because \(N\) is the normal closure of
\(L/K\); equivalently, the action of \(G\) on the conjugate embeddings
of \(L\) is faithful.

For \(Y=\Spec B\), the finite-normalization construction of
Section~\ref{sec:normalized-covers-inertia} gives
\[
  Z=\Norm_N(Y)\longrightarrow
  \overline X=\Norm_L(Y)\longrightarrow Y.
\]
If \(E\in Z^{(1)}\), with decomposition group \(D_E\) and inertia group
\(I_E\), then its conjugate centers are indexed by
\(\mathcal Q_E=D_E\backslash G/H\).  For
\(\omega=D_EgH\), Proposition~\ref{prop:intermediate-ramification-index}
states
\[
  e\bigl(\operatorname{cent}(E,\omega)/\nu(E)\bigr)
  =
  i_E(\omega)
  =
  [I_E:I_E\cap gHg^{-1}].
\]

The following proposition is the standard residual-factor description
of a separable nonnormal cubic extension.  We include its short proof
to fix the marked identification with \(N\).

\begin{proposition}[Cubic residual-factor identification]
\label{prop:cubic-tensor-decomposition}
The polynomial \(h_\alpha\) is irreducible over \(L\), and there is an
\(L\)-algebra isomorphism
\[
  \varphi\colon L[T]/(h_\alpha)\xrightarrow{\sim}_L N.
\]
Consequently,
\[
  \Psi\colon L\otimes_KL\xrightarrow{\sim}_L L\times N,
  \qquad
  \operatorname{pr}_1\circ\Psi=\mu_L.
\]
\end{proposition}

\begin{proof}
Write
\[
  m_\alpha(T)=T^3+a_2T^2+a_1T+a_0.
\]
If the quadratic \(h_\alpha\) were reducible over \(L\), it would have a
root \(\beta\in L\).  Separability gives \(\beta\neq\alpha\), and the
third root would be
\[
  \gamma=-a_2-\alpha-\beta\in L.
\]
Thus \(m_\alpha\) would split over \(L\).  Since
\(L=K(\alpha)\), the extension \(L/K\) would then be the splitting
field of the separable polynomial \(m_\alpha\), hence normal.  This
contradicts the assumed nonnormality, so \(h_\alpha\) is irreducible.

Let \(\beta\) denote the image of \(T\) in \(L[T]/(h_\alpha)\).  This
field contains the three roots
\[
  \alpha,\qquad \beta,\qquad -a_2-\alpha-\beta
\]
of \(m_\alpha\), and is generated over \(L\) by \(\beta\).  It is
therefore the splitting field of \(m_\alpha\), which is the normal
closure \(N\) of \(L/K\).  Choose the resulting \(L\)-algebra
isomorphism
\[
  \varphi\colon L[T]/(h_\alpha)\xrightarrow{\sim}_L N.
\]

Let
\[
  \Phi_\alpha\colon
  L\otimes_KL
  \xrightarrow{\sim}_L
  L\times L[T]/(h_\alpha)
\]
be the marked-root decomposition of
Theorem~\ref{thm:marked-root-tensor-decomposition}.  It satisfies
\[
  \operatorname{pr}_1\circ\Phi_\alpha=\mu_L.
\]
Composing its second factor with \(\varphi\), let
\[
  \Psi=(\id_L\times\varphi)\circ\Phi_\alpha.
\]
Since \(\id_L\times\varphi\) leaves the first factor unchanged,
\(\operatorname{pr}_1\circ\Psi=\mu_L\).
\end{proof}

\begin{theorem}[Galois group of a nonnormal cubic extension]
\label{thm:index-three-core-free-s3}
\[
  \Gal(N/K)\cong S_3.
\]
\end{theorem}

\begin{proof}
Proposition~\ref{prop:cubic-tensor-decomposition} identifies \(N\) with
the degree-two extension \(L[T]/(h_\alpha)\) of \(L\).  Hence
\[
  |H|=[N:L]=2.
\]
The normal-closure identities recalled above give
\[
  \core_G(H)=1,
  \qquad
  [G:H]=[L:K]=3.
\]
Hence the action of \(G\) on the three cosets of \(H\) gives
\[
  \rho\colon G\longrightarrow\operatorname{Perm}(G/H)\cong S_3,
  \qquad
  \ker\rho=\core_G(H)=1.
\]
Core-freeness therefore makes \(\rho\) injective, so
\[
  |G|\leq |S_3|=6.
\]
The index formula gives
\[
  |G|=[G:H]|H|=3\cdot2=6.
\]
Thus \(|G|=6\), and the injective homomorphism \(\rho\) is an
isomorphism \(G\cong S_3\).
\end{proof}

\subsection{The cubic off-diagonal factor}

Retain the cubic data of Section~\ref{sec:cubic-extension}, and suppose
in addition that \(F\) is Keller.  Thus
\[
  B=\C[F_1,F_2,F_3]\subseteq A=\C[x_1,x_2,x_3],
  \qquad
  K=\Frac(B)\subseteq L=\Frac(A)=K(\alpha)\subseteq N,
\]
where \(L/K\) is separable, nonnormal, and of degree three, and \(N/K\)
is its normal closure.  Let
\[
  \theta\colon K\otimes_BA\xrightarrow{\sim}L,
  \qquad
  \lambda\otimes a\longmapsto\lambda a.
\]

\begin{corollary}[Generic cubic \(S_3\) collision]
\label{cor:cubic-s3-collision}
There is a \(K\)-algebra isomorphism
\[
  \Psi_F\colon
  K\otimes_BC_F\xrightarrow{\sim}_K L\times N,
  \qquad
  \operatorname{pr}_1\circ\Psi_F
  =
  \theta\circ(1\otimes\bar\mu_F).
\]
It satisfies
\[
  \Psi_F\!\left(
    \ker\bigl(\theta\circ(1\otimes\bar\mu_F)\bigr)
  \right)
  =
  0\times N,
  \qquad
  \Gal(N/K)\cong S_3.
\]
Consequently,
\[
  \Obs(F)\neq0,
  \qquad
  I_R\subsetneq I_\Delta,
  \qquad
  R_F^\circ\neq\varnothing,
  \qquad
  F\notin\Aut_{\mathrm{poly}}(\A^3_{\C}).
\]
\end{corollary}

\begin{proof}
Theorem~\ref{thm:generic-collision-base-change} gives a \(K\)-algebra
isomorphism
\[
  \Phi_F\colon K\otimes_BC_F\xrightarrow{\sim}_K L\otimes_KL
\]
such that
\[
  \mu_L\circ\Phi_F
  =
  \theta\circ(1\otimes\bar\mu_F).
\]
Proposition~\ref{prop:cubic-tensor-decomposition} gives
\[
  \Psi\colon L\otimes_KL\xrightarrow{\sim}_L L\times N,
  \qquad
  \operatorname{pr}_1\circ\Psi=\mu_L.
\]
After restricting \(\Psi\) from \(L\) to \(K\), let
\[
  \Psi_F=\Psi\circ\Phi_F.
\]
It follows that
\[
  \operatorname{pr}_1\circ\Psi_F
  =
  \theta\circ(1\otimes\bar\mu_F),
  \qquad
  \Psi_F\!\left(
    \ker\bigl(\theta\circ(1\otimes\bar\mu_F)\bigr)
  \right)
  =
  0\times N\neq0.
\]

Suppose that the diagonal-evaluation homomorphism
\(\bar\mu_F\colon C_F\to A\) were injective.  Since \(K\) is a
localization of \(B\), it is flat over
\(B\), so \(1\otimes\bar\mu_F\) would also be injective.
Theorem~\ref{thm:generic-collision-base-change} shows that \(\theta\)
is an isomorphism.  Therefore
\(\theta\circ(1\otimes\bar\mu_F)\) would be injective, contradicting
the nonzero kernel \(0\times N\) above.  Hence
\(\ker(\bar\mu_F)\neq0\), and
Proposition~\ref{prop:obstruction-kernel} gives
\(\Obs(F)\neq0\).  Proposition
\ref{prop:collision-contained-in-diagonal} shows that
\(I_R\subseteq I_\Delta\) strict, and
Corollary~\ref{cor:off-diagonal-empty-criterion} gives
\(R_F^\circ\neq\varnothing\).

Finally, Theorem~\ref{thm:index-three-core-free-s3} identifies
\(\Gal(N/K)\) with \(S_3\).  If \(F\) were a polynomial automorphism,
Theorem~\ref{thm:automorphism-criterion} would give
\(\Obs(F)=0\), contrary to \(\Obs(F)\neq0\).  Since \(F\) is Keller, it is
a dimension-three Jacobian-conjecture counterexample.
\end{proof}

For the same cubic map and field tower, let
\[
  X=\Spec A,
  \qquad
  Y=\Spec B,
  \qquad
  G=\Gal(N/K),
  \qquad
  H=\Gal(N/L).
\]
Suppose there is a diagram over \(Y\)
\[
  \mathcal D=
  \left(
    Z=\Norm_N(Y)\xrightarrow{\pi}
    \overline X=\Norm_L(Y)\xrightarrow{\bar\nu}Y,\,
    j\colon X\hookrightarrow\overline X
  \right)
\]
in which \(\pi\) and \(\bar\nu\) are finite, \(j\) is an open immersion,
\(\bar\nu\circ j=f_F\), and \(f_F\colon X\to Y\) is \'etale.  Let
\[
  \partial X=\overline X\setminus j(X),
  \qquad
  G\cong S_3
\]
by Theorem~\ref{thm:index-three-core-free-s3}.

\begin{corollary}[Hidden inertia in the \(S_3\)-normalization]
\label{cor:cubic-hidden-inertia}
\[
  \Ram^{(1)}(Z/Y)\subseteq\mathcal H(\mathcal D).
\]
\end{corollary}

\begin{proof}
Fix
\[
  E\in\Ram^{(1)}(Z/Y).
\]
By the characteristic-zero description of divisorial ramification,
\(I_E\neq1\).  Because \(N\) is the normal closure of \(L/K\),
\(\core_G(H)=1\).  Proposition~\ref{prop:core-free-sheet-action}
states in this case that
\[
  \exists\omega\in\mathcal Q_E,
  \qquad
  i_E(\omega)>1.
\]
Moreover, Proposition~\ref{prop:visibility-to-boundary}, applied to the
\'etale map \(f_F\), states that for every
\(\omega'\in\mathcal Q_E\),
\[
  i_E(\omega')>1
  \quad\Longrightarrow\quad
  \operatorname{cent}(E,\omega')\in\partial X.
\]
These are exactly the existence and boundary conditions in
Definition~\ref{def:no-hidden-inertia}; hence
\(E\in\mathcal H(\mathcal D)\).  Since \(E\) was arbitrary, the stated
inclusion follows.
\end{proof}

\begin{remark}[Ordered roots]
Let
\[
  m_\alpha(T)=T^3+a_2T^2+a_1T+a_0
\]
have roots \(\alpha_1,\alpha_2,\alpha_3\) in \(N\).  For distinct
\(i,j,k\),
\[
  \alpha_k=-a_2-\alpha_i-\alpha_j,
  \qquad
  K(\alpha_i,\alpha_j)
  =
  K(\alpha_1,\alpha_2,\alpha_3)
  =
  N.
\]
Thus the off-diagonal factor, which marks two distinct roots, is already
the ordered-root normal closure.
\end{remark}

\begin{leanfiles}
\leanpath{CollisionIdeals/ComplexThree/CubicGaloisGroup.lean},
\leanpath{CollisionIdeals/ComplexThree/S3Collision.lean},
\leanpath{CollisionIdeals/ComplexThree/OffDiagonal.lean},
\leanpath{CollisionIdeals/ComplexThree/CubicS3.lean}, and
\leanpath{CollisionIdeals/ComplexThree/JacobianConjecture.lean}.  The
Lean cubic witness takes the residual-field isomorphism and
\(H\neq1\) as explicit inputs; Proposition~\ref{prop:cubic-tensor-decomposition}
derives them here from nonnormality.  Hidden inertia is formalized
pointwise at a supplied ramified divisor.
\end{leanfiles}

\section{Conclusion}
\label{sec:conclusion}

The quotient \(\Obs(F)\) vanishes exactly when equality of images agrees
scheme-theoretically with equality of points.  In dimension two,
\(\mathcal H(\mathcal D_F)=\varnothing\) implies \(\Obs(F)=0\); in the
cubic dimension-three class, its nonzero generic factor is the
ordered-root \(S_3\)-normal closure.  Both statements arise from the same
collision, off-diagonal, and conjugate-sheet constructions.

Under the respective conditions of Sections~\ref{sec:planar-vanishing}
and~\ref{sec:cubic-s3-obstruction}, and with
\(\Ram^{(1)}(Z/Y)\neq\varnothing\) in the cubic normalization, the
contrast is
\[
  \begin{array}{c@{\qquad}|@{\qquad}c}
    \hline
    \text{\rm dimension two}
      & \text{\rm dimension three}\\ \hline
    \text{\rm hidden inertia excluded}
      & \text{\rm hidden inertia in the \(S_3\)-cover realized}\\
    \Obs(F)=0 & \Obs(F)\neq0\\
    I_R=I_\Delta & I_R\subsetneq I_\Delta\\
    F\text{\rm\ is a polynomial automorphism}
      & F\text{\rm\ is not a polynomial automorphism}.\\
    \hline
  \end{array}
\]

\paragraph{Motivating example.}
The explicit three-dimensional polynomial counterexample announced by
Levent Alpöge and independently formalized in Isabelle/HOL by Ramos,
Hulak, and de Queiroz \cite{AFPJacobianCounterexample} motivated the
cubic \(S_3\) comparison.  The formal verification establishes a
constant nonzero Jacobian determinant and distinct rational points with
the same image.  That particular map is not used in the arguments above:
the dimension-three result is the structural implication from the
stated cubic and marked-normal-closure conditions.

The formalization uses Lean~4 and Mathlib \cite{Lean4,Mathlib}.

\bibliographystyle{alpha}
\bibliography{references}

\end{document}
